%% =====================================================================
%%  B3-T-28-01 -- what B3 contributes to "Assume Formlessness"
%%  -------------------------------------------------------------------
%%  LAYER A: APPEND-ONLY. Written once at the entry's write, then left
%%  alone. If the source has more to say later, add a bullet -- do not
%%  rewrite. Material found ONLY here goes in the `unique` slot with
%%  @unique set to yes, which makes it undeletable (rule R10).
%% =====================================================================
%% @kind: source
%% @id: B3-T-28-01
%% @book: B3
%% @topic: T-28-01
%% @locator: Law 48, pp. 419--430
%% @status: distilled
%% @unique: no
%% @integrated: yes
%% @updated: 2026-09-19

\dossier{B3}{T-28-01}{\bookshort{B3}}{Law 48, pp. 419--430}

\begin{distilled}
  \item By taking a shape you give the enemy something to grasp; the protectable position is fluid, adaptable, on the move.
  \item The serpent image: nothing in the animal is fixed --- scale by scale --- and precisely that mobility makes it untouchable.
  \item T. E. Lawrence's desert warfare as the terminal case: no uniform, no fixed position, everywhere and nowhere; the Turks could not hit what never collected into a target.
  \item The closing discipline: nothing is certain and no law is fixed; form is assumed as tactic when it serves and shed when it is priced --- and the law itself warns the advice is not for the decent.
\end{distilled}
